\clearpage\pagestyle{otMain}\fontsize{11}{15}\selectfont\corrigesChapitre{5}{Équations et mise en équation}
\begin{blocLivre}{cor-5-1}\begin{corrige}{1}{Une opération inverse}\demarche{Appliquer l’opération inverse aux deux membres.}
On soustrait 13 à la première égalité : $\res{x=-9}$. On divise la seconde par $-4$ : $\res{x=-9/2}$. Les deux équations sont indépendantes.
\end{corrige}\end{blocLivre}
\begin{blocLivre}{cor-5-2}\begin{corrige}{2}{Deux étapes}\demarche{Isoler le produit avant l’inconnue.}
On ajoute 5 : $7x=35$. Puis on divise par 7 : $\res{x=5}$. Contrôle : $7\times5-5=30$.
\end{corrige}\end{blocLivre}
\begin{blocLivre}{cor-5-3}\begin{corrige}{3}{L’inconnue des deux côtés}\demarche{Rassembler les termes inconnus avant les constantes.}
On soustrait $5x$ : $3x+3=-12$. On retire 3 : $3x=-15$. Ainsi $\res{x=-5}$. Les deux membres initiaux valent $-37$.
\end{corrige}\end{blocLivre}
\begin{blocLivre}{cor-5-4}\begin{corrige}{4}{Avec une parenthèse}\demarche{Développer avant de regrouper.}
On développe : $6x-15=4x+9$, donc $2x=24$ et $\res{x=12}$. Vérification : $3(24-5)=57=48+9$.
\end{corrige}\end{blocLivre}
\begin{blocLivre}{cor-5-5}\begin{corrige}{5}{Une solution fractionnaire}\demarche{Conserver une valeur exacte lors de la division.}
L’égalité donne $5x=7$, puis $\res{x=7/5}$. On peut écrire 1,4 ici car le décimal est exact. Une fraction n’est pas une erreur de résolution.
\end{corrige}\end{blocLivre}
\begin{blocLivre}{cor-5-6}\begin{corrige}{6}{Une équation fractionnaire}\demarche{Multiplier chaque terme par le même dénominateur commun.}
On multiplie toute l’égalité par 12 : $4x+3x=168$. Donc $7x=168$ et $\res{x=24}$. Contrôle : $24/3+24/4=8+6=14$.
\end{corrige}\end{blocLivre}
\begin{blocLivre}{cor-5-7}\begin{corrige}{7}{Tout ou rien}\demarche{Interpréter l’égalité restante lorsque l’inconnue disparaît.}
La première égalité se réduit à $4x-8=4x-8$ : tous les réels sont solutions. La seconde imposerait $-8=1$ après soustraction de $4x$ : elle n’a aucune solution.
\end{corrige}\end{blocLivre}
\begin{blocLivre}{cor-5-8}\begin{corrige}{8}{Un produit nul}\demarche{Séparer les deux facteurs, pas les multiplier inutilement.}
Au moins un facteur est nul : $2x-7=0$ ou $x+4=0$. On obtient $\res{x=7/2\text{ ou }x=-4}$. Les deux valeurs annulent bien le produit.
\end{corrige}\end{blocLivre}
\begin{blocLivre}{cor-5-9}\begin{corrige}{9}{Une équation de carré}\demarche{Examiner le signe du second membre et conserver les deux racines éventuelles.}
Les solutions sont respectivement $\res{-7\text{ et }7}$, $\res{-2\sqrt3\text{ et }2\sqrt3}$, puis aucune solution réelle. Un carré réel n’est jamais négatif.
\end{corrige}\end{blocLivre}
\begin{blocLivre}{cor-5-10}\begin{corrige}{10}{Ne pas perdre zéro}\demarche{Factoriser avant une division par une inconnue.}
On factorise $x(x-6)=0$, donc $\res{x=0\text{ ou }x=6}$. Diviser par $x$ exigerait $x\ne0$ et supprimerait la première solution. Le produit nul conserve toutes les possibilités.
\end{corrige}\end{blocLivre}
\begin{blocLivre}{cor-5-11}\begin{corrige}{11}{Des billets de spectacle}\demarche{Exprimer les deux effectifs avec une seule inconnue.}
Notons $x$ le nombre de jeunes ; il y a $28-x$ autres billets. L’équation est $6x+10(28-x)=232$, soit $280-4x=232$. Ainsi $\res{x=12}$. Il reste 16 autres billets ; $72+160=232$. La solution est entière et comprise entre 0 et 28.
\end{corrige}\end{blocLivre}
\begin{blocLivre}{cor-5-12}\begin{corrige}{12}{Un rectangle à retrouver}\demarche{Traduire le périmètre avec les deux côtés.}
Notons $x$ sa largeur, sa longueur vaut $x+7$. Alors $2x+2(x+7)=58$, donc $4x=44$ et $x=11$. Les dimensions sont $\res{11\text{ cm et }18\text{ cm}}$, l’aire $\res{198\text{ cm}^2}$.
\label{fig-cor-5-12}\figureGeo[corrige]{rectangle-equation}
\end{corrige}\end{blocLivre}
\begin{blocLivre}{cor-5-13}\begin{corrige}{13}{Retrouver le nombre initial}\demarche{Nommer le nombre initial tout au long du programme.}
Pour un départ $x$, le résultat vaut $3(x+4)-2x=x+12$. L’équation $x+12=25$ donne $\res{x=13}$. Vérification : $17\times3-26=25$. La traduction distingue le résultat courant et le nombre initial.
\end{corrige}\end{blocLivre}
\begin{blocLivre}{cor-5-14}\begin{corrige}{14}{Un achat rentable}\demarche{Calculer le gain annuel puis tester les entiers autour du seuil.}
Après $n$ années, les dépenses sont $38+12n$ et $30n$. Il faut économiser au moins les 38 € initiaux ; le gain annuel est $30-12=18$ €. Or $2\times18=36<38$ et $3\times18=54\ge38$. La première année entière est $\res3$. Contrôle : 74 € contre 90 € à trois ans.
\end{corrige}\end{blocLivre}
\begin{blocLivre}{cor-5-15}\begin{corrige}{15}{Deux aires égales}\demarche{Développer pour faire disparaître les termes carrés communs.}
On pose $(x+2)^2=x(x+5)$. Après développement : $x^2+4x+4=x^2+5x$, donc $\res{x=4}$. Le carré a pour côté 6 et le rectangle pour côtés 4 et 9. L’aire commune est $\res{36\text{ cm}^2}$. Toutes les dimensions sont positives.
\label{fig-cor-5-15}\figureGeo[corrige]{aires-egales}
\end{corrige}\end{blocLivre}
